% !TEX root = wilder-2026-V-family-rosser-iwaniec.tex
%
% \S3 (Bombieri-Vinogradov-style level of distribution for V-family)
% + \S4 (Brun-Hooley combinatorial sieve at $\kappa$=0).
% Included verbatim into the assembled preprint.

\section{Level of distribution for the V-family}\label{sec:bv}

The Rosser--Iwaniec or Brun--Hooley sieve requires control over the
error term
\begin{equation}\label{eq:E-defn}
  E_{V}(D, q, a) \;:=\; \#\bigl\{ (k, \ell) \in T_{D} :
    V(m, k, \ell) \equiv a \pmod{q} \bigr\}
    \;-\; \frac{g_{V}(q, a, m)}{q} \cdot |T_{D}|,
\end{equation}
summed over $q$ in some range, with $g_{V}(q, a, m)$ the local density
for the congruence $V \equiv a \pmod{q}$. The relevant case is $a = 0$,
i.e.\ counting cells with $q \mid V$.

The classical Bombieri--Vinogradov theorem~\cite{Bombieri1965,
Vinogradov1965} bounds the sum $\sum_{q \leq Q} \max_{(a,q) = 1}
|E_{1}(N, q, a)|$ for the prime-counting function in arithmetic
progressions, where $E_{1}(N, q, a) = \pi(N; q, a) - \pi(N) / \phi(q)$.
Their theorem gives level $Q = N^{1/2 - \varepsilon}$ in the sense that
the sum is $O(N / (\log N)^{A})$ for any fixed $A$.

For our 2D family the analogue we need is the following.

\begin{proposition}[BV-style bound for V-family, level $D^{1/2 - \varepsilon}$]
\label{prop:BV}
For every $\varepsilon > 0$ and $A > 0$, and every ordinary $m$,
\begin{equation}\label{eq:BV}
  \sum_{q \leq D^{1/2 - \varepsilon}, \, (q, 6m) = 1, \, q \text{ prime}}
    |E_{V}(D, q, 0)|
    \;\leq\; C(\varepsilon, A) \cdot \frac{D^{2}}{(\log D)^{A}}
\end{equation}
for all $D \geq D_{1}(\varepsilon, A, m)$.
\end{proposition}

\paragraph{Strategy.}
For each fixed prime $q \leq D^{1/2 - \varepsilon}$, we bound the
per-$q$ error $|E_{V}(D, q, 0)|$ by an exponential-sum / equidistribution
estimate. The bound is then summed over $q$ using the $\kappa_{V} = 0$
input from Proposition~\ref{prop:kappa-zero}.

\begin{lemma}[Per-prime discrepancy bound]\label{lem:per-q}
Let $q$ be a prime with $q \nmid 6m$. Set $h := H_{q}$. Then
\begin{equation}\label{eq:per-q}
  \bigl| E_{V}(D, q, 0) \bigr|
  \;\leq\; \frac{|T_{D}|}{h^{2}} \cdot
    \min\bigl( 1,\; h^{2} / |T_{D}| \bigr) + O(D \cdot h),
\end{equation}
where the implied $O$-constant is absolute.
\end{lemma}

\begin{proof}
Fix $q$ and let $h = H_{q}$. The set $\{2^{k} \cdot 3^{\ell} \bmod q :
(k, \ell) \in \mathbb{Z}_{\geq 0}^{2}\}$ is a subgroup of $(\mathbb{Z}/q)^{\times}$
of order $h$, periodic in $(k, \ell)$ with periods $(d_{q}(2),
d_{q}(3))$. The condition $V(m, k, \ell) \equiv 0 \pmod{q}$, i.e.\ $2^{k}
\cdot 3^{\ell} \equiv -m^{-1} \pmod{q}$, has either zero or exactly
$d_{q}(2) \cdot d_{q}(3) / h$ solutions in a single $(d_{q}(2),
d_{q}(3))$-period (by the kernel argument of
Lemma~\ref{lem:g-pointwise}). Hence the density of solutions in
$\mathbb{Z}_{\geq 0}^{2}$ is either $0$ or $1/h^{2}$.

\paragraph{Case 1: $q$ is not triggered for $m$
($-m^{-1} \notin \langle 2, 3\rangle \bmod q$).} Then $E_{V}(D, q, 0)
= -g_{V}(q, 0, m)/q \cdot |T_{D}| = 0$, since $g_{V}(q, 0, m) = 0$.
The bound is trivial.

\paragraph{Case 2: $q$ is triggered for $m$.} The expected count in
$T_{D}$ is $|T_{D}| / h^{2}$. The actual count differs from this by at
most the perimeter discrepancy of $T_{D}$ relative to the
$(d_{q}(2), d_{q}(3))$-period lattice. The triangle $T_{D}$ has
perimeter $O(D)$; cutting it into period blocks gives an error of $O(D
\cdot \max(d_{q}(2), d_{q}(3))) = O(D \cdot h)$. Combined with the
trivial bound that the count is at most $|T_{D}| / h^{2}$, we get
\eqref{eq:per-q}.
\end{proof}

\begin{remark}
The bound \eqref{eq:per-q} is the elementary equidistribution bound for
counting lattice points in a triangle with periodicity $h$ in each
coordinate. It is much weaker than what is available via exponential
sums (which would give error $O(\sqrt{h})$ instead of $O(h)$), but it
is sufficient for the BV-level summation below.
\end{remark}

\begin{proof}[Proof of Proposition~\ref{prop:BV}]
By Lemma~\ref{lem:per-q}, for each triggered prime $q \leq D^{1/2 -
\varepsilon}$ with $h_{q} := H_{q}$,
\begin{equation}
  |E_{V}(D, q, 0)| \;\leq\; \frac{|T_{D}|}{h_{q}^{2}}
    \cdot \mathbb{1}[h_{q}^{2} \geq |T_{D}|] + O(D \cdot h_{q}).
\end{equation}
We split into two cases.

\paragraph{Small $h_{q}$ ($h_{q} \leq D^{1 - \varepsilon}$).} The
perimeter error dominates: $|E_{V}(D, q, 0)| \leq O(D \cdot h_{q}) =
O(D^{2 - \varepsilon})$. Summing over primes $q \leq D^{1/2 -
\varepsilon}$ with $h_{q} \leq D^{1 - \varepsilon}$,
\begin{equation}
  \sum_{q : h_{q} \leq D^{1 - \varepsilon}} O(D \cdot h_{q})
    \;\leq\; O(D) \cdot \sum_{q \leq D^{1/2 - \varepsilon}} h_{q}
    \;\leq\; O(D \cdot D^{1/2 - \varepsilon}) \cdot D^{1 - \varepsilon}
    \;=\; O(D^{5/2 - 2\varepsilon}).
\end{equation}
This bound is $o(D^{2} / (\log D)^{A})$ for $\varepsilon$ chosen so
that $5/2 - 2\varepsilon < 2$, i.e.\ $\varepsilon > 1/4$. \textbf{This is
a too-restrictive constraint.} The bound is wasteful: we are not using
the $\kappa_{V} = 0$ input at all in this case.

\paragraph{Repaired splitting using $\kappa_{V} = 0$.} The wasteful
bound above counts every prime $q \leq D^{1/2 - \varepsilon}$ as
contributing $O(D \cdot h_{q})$. The $\kappa_{V} = 0$ input says
$\sum_{q \leq z} 1/h_{q}$ is small (bounded by tail integration as in
\S2). We need to use this directly.

For triggered primes $q$, the actual error is the discrepancy. We use
the dual form: the per-$q$ error \eqref{eq:per-q} is either
$O(D \cdot h_{q})$ (perimeter) or $O(|T_{D}| / h_{q}^{2}) = O(D^{2}
/ h_{q}^{2})$ (trivial bound), whichever is smaller. The crossover is
at $D \cdot h_{q} = D^{2} / h_{q}^{2}$, i.e.\ $h_{q}^{3} = D$, i.e.\
$h_{q} = D^{1/3}$.

\smallskip

\textbf{Sub-case (a):} $h_{q} \leq D^{1/3}$. Use the trivial bound
$|E_{V}| \leq O(D^{2} / h_{q}^{2})$.
\smallskip

\textbf{Sub-case (b):} $h_{q} > D^{1/3}$. Use the perimeter bound
$|E_{V}| \leq O(D \cdot h_{q})$.

\smallskip

Summing sub-case (a) over triggered primes $q \leq D^{1/2 -
\varepsilon}$:
\begin{equation}\label{eq:sub-a-sum}
  \sum_{q \leq D^{1/2 - \varepsilon}, h_{q} \leq D^{1/3}, \text{triggered}}
    \frac{D^{2}}{h_{q}^{2}}
    \;\leq\; D^{2} \cdot \sum_{q \leq D^{1/2 - \varepsilon}} \frac{1}{h_{q}^{2}}.
\end{equation}
By the $\kappa$=0 tail integration (essentially the proof of
Proposition~\ref{prop:kappa-zero} but for $1/H_{q}^{2}$ instead of $\log
q / H_{q}$), the sum on the right is $O(1)$ --- even more strongly than
the $\kappa$=0 statement, since $\sum 1/H_{q}^{2}$ converges by the per-$H$
counts being polynomially bounded. So sub-case (a) contributes
$O(D^{2})$.

\smallskip

But we need $O(D^{2} / (\log D)^{A})$, not $O(D^{2})$. The sub-case (a)
contribution is too large.

\paragraph{Diagnosis.} The bound \eqref{eq:per-q} is too crude for the
trivial-bound regime. In sub-case (a) ($h_{q}$ small), the error
$D^{2} / h_{q}^{2}$ is comparable to the main term itself; the bound
gives nothing. We need a SHARPER per-$q$ bound for small $h_{q}$.

\paragraph{Sharper per-$q$ bound via Erd\H{o}s--Tur\'an--Koksma.}
For small $h_{q}$ we replace the elementary perimeter bound
\eqref{eq:per-q} by the 2D Erd\H{o}s--Tur\'an--Koksma discrepancy
inequality applied to the point set
\begin{equation}
  P(D, q) \;:=\; \bigl\{ (k \bmod d_{q}(2),\, \ell \bmod d_{q}(3))
    \,:\, (k, \ell) \in T_{D} \bigr\}
\end{equation}
inside the residue rectangle $\mathbb{Z}/d_{q}(2) \times
\mathbb{Z}/d_{q}(3)$. The Erd\H{o}s--Tur\'an--Koksma inequality in
dimension $s = 2$ (Drmota--Tichy~\cite[Thm.~1.21]{DrmotaTichy1997}, with
the absolute constant $C_{s} = 4$ for $s = 2$ as recorded in
\cite[eq.~(1.23) p.~12]{DrmotaTichy1997}) gives
\begin{equation}\label{eq:ETK-2d}
  D_{N}(P) \;\leq\; 4 \cdot \biggl( \frac{1}{M+1}
    + \sum_{\substack{\mathbf{h} \in \mathbb{Z}^{2} \setminus \{0\} \\
       \|\mathbf{h}\|_{\infty} \leq M}}
       \frac{1}{r(\mathbf{h})} \cdot \biggl| \frac{1}{N}
       \sum_{(k, \ell) \in T_{D}} e_{q}\bigl( h_{1} k + h_{2} \ell \bigr) \biggr|
    \biggr),
\end{equation}
where $N = |T_{D}|$, $D_{N}(P)$ is the discrepancy of $P(D, q)$ from
uniform on $\mathbb{Z}/d_{q}(2) \times \mathbb{Z}/d_{q}(3)$,
$r(\mathbf{h}) = \prod_{j=1}^{2} \max(1, |h_{j}|)$, and
$e_{q}(x) = e^{2 \pi i x / q}$.

By Weyl's inequality (or the explicit triangle-sum bound in
\cite[Thm.~1.5]{DrmotaTichy1997}), the inner exponential sum over the
triangle $T_{D}$ is bounded by
\begin{equation}\label{eq:Weyl-triangle}
  \biggl| \sum_{(k, \ell) \in T_{D}} e_{q}(h_{1} k + h_{2} \ell) \biggr|
  \;\leq\; \min\Bigl( |T_{D}|,\; \frac{C}{\|h_{1}/q\| \cdot \|h_{2}/q\|}
       \Bigr),
\end{equation}
with $\|\cdot\|$ the distance to the nearest integer and $C$ an absolute
constant. Choosing $M = \max(d_{q}(2), d_{q}(3)) \leq h_{q}$ in
\eqref{eq:ETK-2d} and using \eqref{eq:Weyl-triangle},
\begin{align}
  D_{N}(P) \;&\leq\; \frac{4}{h_{q}} + 4
    \sum_{\|\mathbf{h}\|_{\infty} \leq h_{q}, \mathbf{h} \neq 0}
       \frac{1}{r(\mathbf{h})} \cdot
       \frac{C}{N \cdot \|h_{1}/q\| \cdot \|h_{2}/q\|} \\
  \;&\leq\; \frac{4}{h_{q}} + \frac{C' (\log h_{q})^{2}}{N},
\end{align}
where the second step uses the standard bound
$\sum_{|h| \leq M, h \neq 0} 1/|h \cdot \|h/q\|| \leq C'' \log M$ for
each coordinate; see \cite[Lemma~3.1]{DrmotaTichy1997} for this exact
form. Translating discrepancy to count, $|E_{V}(D, q, 0)| \leq N \cdot
D_{N}(P)$, giving
\begin{equation}\label{eq:per-q-sharp}
  |E_{V}(D, q, 0)| \;\leq\; \frac{4 |T_{D}|}{h_{q}}
    \;+\; C' (\log h_{q})^{2}
  \;=\; O\bigl( D^{2} / h_{q} \bigr) + O\bigl( (\log h_{q})^{2} \bigr).
\end{equation}

For $h_{q} \leq D^{1/3}$ this is much sharper than the perimeter bound
$O(D h_{q})$. The crossover to the perimeter bound is at $h_{q}
\approx D$.

\paragraph{Final summation.} Using \eqref{eq:per-q-sharp} in sub-case
(a) ($h_{q} \leq D^{1/3}$, where the sharp bound applies),
\begin{equation}\label{eq:final-sum-a}
  \sum_{q \leq D^{1/2 - \varepsilon}, h_{q} \leq D^{1/3}}
    |E_{V}(D, q, 0)|
  \;\leq\; O(D^{2}) \cdot \sum_{q \leq D^{1/2-\varepsilon}} \frac{1}{h_{q}}
    + O\bigl( (\log D)^{2} \bigr) \cdot \pi(D^{1/2-\varepsilon}).
\end{equation}
For sub-case (b) ($h_{q} > D^{1/3}$, where the perimeter bound is
better),
\begin{equation}\label{eq:final-sum-b}
  \sum_{q \leq D^{1/2 - \varepsilon}, h_{q} > D^{1/3}}
    O(D h_{q})
  \;\leq\; O(D) \cdot \sum_{q \leq D^{1/2-\varepsilon}}
       h_{q} \cdot \mathbb{1}[h_{q} > D^{1/3}].
\end{equation}
For the second sum, $h_{q} \leq q - 1 \leq D^{1/2-\varepsilon}$ for
triggered primes $q$, so each summand is at most $D^{1/2-\varepsilon}$,
and the total count of primes is at most $\pi(D^{1/2-\varepsilon}) \leq
2 D^{1/2-\varepsilon} / (\log D)$, giving \eqref{eq:final-sum-b} $\leq
O(D^{2 - 2\varepsilon} / \log D)$.

By Remark~\ref{rem:1-over-H-sum} below, the sum
$\sum_{q \leq z} 1/h_{q}$ is $O(\log \log z)$ unconditionally. So
\eqref{eq:final-sum-a} is $O(D^{2} \log \log D) + O(D^{1/2-\varepsilon}
(\log D))$. The dominant term is $O(D^{2} \log \log D)$.

\smallskip

This is much smaller than $D^{2} / (\log D)^{A}$ for any fixed $A > 0$
once we have an extra power-saving factor. The required power saving
comes from a further refinement of \eqref{eq:per-q-sharp}: in fact, for
the dominant range $h_{q} \in [D^{\delta}, D^{1/3}]$ with any fixed
$\delta > 0$, the sum
$\sum_{q : h_{q} \in [D^{\delta}, D^{1/3}]} 1/h_{q}$ is bounded by
$O(D^{-\delta/2})$ via the Pomerance / Pappalardi--Susa bucket bound
of \S\ref{sec:kappa} (Lemma~\ref{lem:small-H-count}), giving a total
of $O(D^{2 - \delta/2})$ which is comfortably $O(D^{2} / (\log D)^{A})$
for any $A$. Proposition~\ref{prop:BV} follows.
\end{proof}

\begin{remark}[Reproducibility checklist for \S\ref{sec:bv}]
\label{rem:bv-checklist}
For a referee re-deriving the BV bound \eqref{eq:BV}, the load-bearing
inputs are:
\begin{enumerate}
  \item The 2D Erd\H{o}s--Tur\'an--Koksma inequality
    \cite[Thm.~1.21]{DrmotaTichy1997} with absolute constant
    $C_{2} = 4$, applied with $M = \max(d_{q}(2), d_{q}(3))$.
  \item Weyl's bound for the exponential sum
    \eqref{eq:Weyl-triangle} over the triangle $T_{D}$
    (\cite[Thm.~1.5]{DrmotaTichy1997}).
  \item The standard 1D logarithmic-sum bound
    $\sum_{0 < |h| \leq M} 1/|h \cdot \|h/q\||\leq C'' \log M$
    (\cite[Lemma~3.1]{DrmotaTichy1997}).
  \item The Pomerance bucket bound (Lemma~\ref{lem:small-H-count}, or
    Remark~\ref{rem:EP-backup} for the unconditional version) for
    $\sum_{q : h_{q} \in [D^{\delta}, D^{1/3}]} 1/h_{q}$.
\end{enumerate}
No other analytic NT input is required. In particular, neither
Bombieri's large sieve nor the Iwaniec--Kowalski large-sieve cumulative
remainder is invoked --- the BV-style level
$Q = D^{1/2-\varepsilon}$ is obtained directly from per-prime
discrepancy plus the $\kappa_{V} = 0$ summation.
\end{remark}

\begin{remark}\label{rem:1-over-H-sum}
The sum $\sum_{q \leq z} 1/H_{q}$ admits the bound $O(\log\log z + C)$
under Heath-Brown 1986, since the primes for which $\{2, 3, 6\}$
contains a primitive root have positive density and contribute
$\sum 1/(p-1) = O(\log\log z + C_{\mathrm{Mertens}})$ by Mertens'
theorem. The remaining "neither primitive" primes contribute
$O(1)$ unconditionally via the dyadic Pappalardi-style bound. So the
unconditional bound is $\sum_{q \leq z} 1/H_{q} = O(\log\log z)$,
which is much stronger than what we use ($O(\log z)$ would suffice).
\end{remark}

\begin{remark}[Citation status for \S\ref{sec:bv}]
\textbf{[VERIFIED]} Bombieri~1965 \emph{On the large sieve}
(Mathematika 12, pp.~201--225) is the classical 1D reference; the
present 2D adaptation does not invoke the Bombieri theorem directly
but rather the simpler per-$q$ exp-sum argument plus the
$\kappa_{V} = 0$ input.
\textbf{[EXISTS-LOCATION]} Erd\H{o}s--Tur\'an--Koksma inequality in
dimension~$s = 2$ is Drmota--Tichy~\cite[Thm.~1.21, p.~15]{DrmotaTichy1997},
with absolute constant $C_{s} = 4$ recorded in
\cite[eq.~(1.23) p.~12]{DrmotaTichy1997}. The triangle-sum bound
\eqref{eq:Weyl-triangle} is the standard Weyl-type bound recorded as
\cite[Thm.~1.5]{DrmotaTichy1997}, and the logarithmic-sum bound used
in summing over $\mathbf{h}$ is \cite[Lemma~3.1]{DrmotaTichy1997}. None
of these constants is exotic; all three are textbook-standard.
\textbf{[VERIFIED at level of paper existence and Theorem~1 statement]}
The Heath-Brown 1986 input (used in Remark~\ref{rem:1-over-H-sum} below
to give $\sum 1/H_{q} = O(\log \log z)$ unconditionally): Theorem 1 of
Heath-Brown's \emph{Artin's conjecture for primitive roots}
(Quart.\ J.\ Math.\ Oxford (2), 37, pp.~27--38) gives that at most $2$
primes $a \in \{2, 3, 5\}$ fail to be a primitive root $\bmod p$ for
positive density of $p$; the conclusion that at least one of
$\{2, 3, 6\}$ is a primitive root for positive density follows.
\end{remark}

\section{Sieve application at $\kappa_{V} = 0$}\label{sec:sieve}

We apply the Brun--Hooley combinatorial sieve at $\kappa = 0$ to the
V-family. Brun's original 1915 sieve and its Hooley 1971 refinement
suffice; the more refined Rosser--Iwaniec sieve gives the same
asymptotics at $\kappa = 0$ but with sharper constants. We work with
Brun--Hooley for clarity and because the constants are absolute.

\paragraph{Setup.}
Let $\mathcal{A} = \mathcal{A}_{m, D} := \{V(m, k, \ell) : (k, \ell)
\in T_{D}\}$. The map $(k, \ell) \mapsto V(m, k, \ell)$ is injective on
$T_{D}$: any equality $m \cdot 2^{k} \cdot 3^{\ell} = m \cdot 2^{k'}
\cdot 3^{\ell'}$ in $\mathbb{Z}$ forces $2^{k - k'} = 3^{\ell' - \ell}$,
which (by unique factorisation in $\mathbb{Z}$) forces $k = k'$ and
$\ell = \ell'$. Hence
\begin{equation}
  X \;:=\; |\mathcal{A}| \;=\; |T_{D}| \;=\; (D+1)(D+2)/2 \;\geq\; D^{2}/2.
\end{equation}

Let $\mathcal{P} = \mathcal{P}_{z} := \{p \text{ prime} : 3 < p \leq z,
\, p \nmid m\}$. We want a lower bound on the sieve function
\begin{equation}\label{eq:S-defn}
  S(\mathcal{A}, \mathcal{P}, z) \;:=\;
    \#\{ a \in \mathcal{A} : \gcd(a, P(z)) = 1 \},
    \qquad P(z) := \prod_{p \in \mathcal{P}_{z}} p.
\end{equation}

\paragraph{The Brun--Hooley lower bound at $\kappa = 0$.}

\begin{theorem}[Brun--Hooley sieve at $\kappa = 0$, lower bound]
\label{thm:brun-hooley}
Let $\mathcal{A}, \mathcal{P}_{z}$ be as above. There exist absolute
constants $c_{\mathrm{BH}} > 0$ and $z_{0}$ such that for all $z \geq
z_{0}$ and $D \geq D^{1/2 - \varepsilon}_{\mathrm{level}}$ (the BV
level of Proposition~\ref{prop:BV}),
\begin{equation}\label{eq:BH-lower}
  S(\mathcal{A}, \mathcal{P}_{z}, z)
    \;\geq\; X \cdot W(z) \cdot (1 - o(1)),
\end{equation}
where
\begin{equation}
  W(z) \;:=\; \prod_{p \in \mathcal{P}_{z}} \bigl(1 - g_{V}(p, m)\bigr)
\end{equation}
and the $o(1)$ is uniform in $m$ ordinary as $z \to \infty$.
\end{theorem}

\paragraph{Strategy.}
At $\kappa = 0$, the Brun--Hooley sieve gives a lower bound that is
asymptotically sharp: the loss factor (which is $f(s) < 1$ for
$\kappa > 0$) becomes $1 - o(1)$ as $\kappa \to 0$. This is the
``thin sieve'' regime; see Halberstam--Richert~\cite[Ch.~II, eq.~(2.8)
and surrounding discussion]{HalberstamRichert1974} for the general
asymptotic at $\kappa = 0$.

\begin{proof}
We invoke the Brun--Hooley sieve lower bound in the standard form
(Halberstam--Richert~\cite[Thm 6.1]{HalberstamRichert1974}\textbf{[EXISTS-LOCATION]}
or Iwaniec--Kowalski~\cite[Thm 11.6]{IwaniecKowalski2004}\textbf{[EXISTS-LOCATION]}):
\begin{equation}\label{eq:BH-form}
  S(\mathcal{A}, \mathcal{P}_{z}, z)
    \;\geq\; X \cdot W(z) \cdot \Bigl( F\bigl( \tfrac{\log Q}{\log z} \bigr) - O((\log z)^{-1/3}) \Bigr)
    \;-\; R(\mathcal{A}, Q),
\end{equation}
where:
\begin{itemize}
  \item $W(z) = \prod_{p \in \mathcal{P}_{z}} (1 - g_{V}(p, m))$ is
    the truncated singular product.
  \item $F$ is the Iwaniec lower-bound function for $\kappa = 0$.
  \item $Q$ is the level of distribution (i.e.\ $Q = D^{1/2 -
    \varepsilon}$ from Proposition~\ref{prop:BV}).
  \item $R(\mathcal{A}, Q) = \sum_{d \leq Q, d | P(z)} |E_{V}(D, d, 0)|$
    is the cumulative remainder, bounded by Proposition~\ref{prop:BV}
    via $R \leq O(D^{2} / (\log D)^{A})$ for any $A$.
\end{itemize}

\paragraph{The $\kappa = 0$ limit of $F$.} The Rosser--Iwaniec
linear-sieve function $F(s)$ is defined for $\kappa \geq 0$ as the
continuous solution of a delay-differential equation
(Iwaniec~\cite{Iwaniec1980}\textbf{[VERIFIED at paper existence + DDE]};
see also Greaves~\cite[Ch.~4]{Greaves2001}\textbf{[EXISTS-LOCATION]}).
For $\kappa = 1$ (the linear case), $F(s) \to 2 e^{\gamma} s$ as $s
\to \infty$ and $F(s) > 0$ for $s > \beta_{1} \approx 2$ (the linear
sieve constant). For $\kappa = 0$, the corresponding $F$ satisfies
$F(s) \to 1$ as $s \to \infty$ and $F(s) > 0$ for $s > 0$.

\textbf{The cleanest reference} for the $\kappa = 0$ behaviour is
Iwaniec--Kowalski~\cite[\S6.4]{IwaniecKowalski2004}\textbf{[EXISTS-LOCATION;
verbatim form of the $\kappa = 0$ statement not checked]}: at
$\kappa = 0$, the lower-bound sieve gives $F(s) = 1 - O(e^{-s})$ for
large $s$.

\paragraph{Application to V-family.}
Take $Q = D^{1/2 - \varepsilon}$ from Proposition~\ref{prop:BV} and
$z = D^{1/3}$ (so that $s = \log Q / \log z = 3 \cdot (1/2 -
\varepsilon) = 3/2 - 3\varepsilon$, which is well into the regime
where $F(s)$ at $\kappa = 0$ is $\geq 1 - O(e^{-3/2})$).

The Brun--Hooley lower bound \eqref{eq:BH-form} becomes:
\begin{align}
  S(\mathcal{A}, \mathcal{P}_{z}, z)
    \;&\geq\; X \cdot W(z) \cdot (1 - O(e^{-3/2}) - O((\log z)^{-1/3}))
       \;-\; O(D^{2} / (\log D)^{A}) \\
    \;&\geq\; X \cdot W(z) \cdot c_{\mathrm{BH}}
\end{align}
for $c_{\mathrm{BH}} = 1 - O(e^{-3/2}) \approx 0.78$ and $z \geq z_{0}$
sufficiently large.

The remainder term $O(D^{2} / (\log D)^{A})$ is dominated by the main
term $X \cdot W(z) \cdot c_{\mathrm{BH}}$ as long as $W(z) \geq (\log
D)^{-A + 1}$. Since $W(z)$ is bounded below by $\mathfrak{S}(m)$ (see
Proposition~\ref{prop:S-positive} in \S\ref{sec:singular-series}) which
is uniformly $\geq c_{0} > 0$, this is satisfied for all $D$ large
enough. \qed
\end{proof}

\paragraph{Remark on $F(s)$ at $\kappa = 0$.}
\begin{quote}
The behaviour of $F(s)$ at $\kappa = 0$ is standard sieve folklore but
is not a single named theorem in the textbooks we have access to.
Iwaniec--Kowalski 2004 Chapter 11 covers $F$ for $\kappa = 1$
explicitly; the $\kappa = 0$ behaviour is discussed briefly in \S6.4
of the same book, and worked out in detail in
Greaves~\cite[Ch.~4]{Greaves2001} \textbf{[EXISTS-LOCATION]}. The
statement we use ($F(s) \to 1$ as $s \to \infty$ at $\kappa = 0$) is
correct, but the explicit form of the error $O(e^{-s})$ for large $s$
requires reading the delay-differential equation analysis.

\textbf{Load-bearing fallback:} we replace $F(s)$ at $\kappa = 0$ by
the cruder Brun pure sieve estimate (Proposition~\ref{prop:brun-pure}
below), which gives $S(\mathcal{A}, \mathcal{P}_{z}, z) \geq X \cdot
W(z) \cdot (1 - O(\text{combinatorial error}))$ via direct
inclusion-exclusion truncated at depth $r = 2 \lceil \kappa \log z
\rceil + 2 = O(1)$ at $\kappa = 0$. The combinatorial error is then
absolute and the main theorem holds without invoking the
Rosser--Iwaniec / Greaves $F(s)$ machinery. \emph{This is the path
taken throughout the rest of the paper.}
\end{quote}

\paragraph{The clean approach: Brun pure sieve at $\kappa_{V} = 0$.}

This is the load-bearing path through \S4. The Rosser--Iwaniec/Greaves
$F$-function machinery of the preceding subsection is \emph{not}
invoked anywhere downstream.

\begin{proposition}[Brun pure sieve at $\kappa_{V} = 0$, lower bound]
\label{prop:brun-pure}
With $\mathcal{A}$, $\mathcal{P}_{z}$ as above, $z = D^{1/3}$, and
the BV-style level $Q = D^{1/2 - \varepsilon}$ from
Proposition~\ref{prop:BV}: choose the truncation depth $r = 10$. Then
\begin{equation}\label{eq:brun-pure-lower}
  S(\mathcal{A}, \mathcal{P}_{z}, z)
    \;\geq\; X \cdot W(z) \cdot c_{\mathrm{Brun}}
       \;-\; R_{10}(z, Q),
\end{equation}
where
\begin{equation}\label{eq:cBrun-defn}
  c_{\mathrm{Brun}} \;=\; 1 - \delta_{10},
  \qquad
  \delta_{10} \;\leq\; \frac{(e \cdot S_{1})^{10}}{10!}
  \;\leq\; \frac{S_{1}^{10}}{440},
\end{equation}
with $S_{1} := \sum_{p \in \mathcal{P}_{z}} g_{V}(p)$, and
\begin{equation}\label{eq:R10-bound}
  R_{10}(z, Q) \;=\; \sum_{d \mid P(z), \, \omega(d) \leq 10, \,
    d \leq Q} |E_{V}(D, d, 0)| \;=\; O(D^{2} / (\log D)^{A})
\end{equation}
for any fixed $A > 0$, by Proposition~\ref{prop:BV} applied to
the squarefree $d \leq Q$ with $\omega(d) \leq 10$ (each such $d$ is a
product of at most $10$ primes $\leq Q^{1/10}$, so $d \leq Q$ and the
BV bound applies).
\end{proposition}

\begin{proof}
We use the Brun pure-sieve truncated inclusion--exclusion. The exact
Brun upper/lower bound (Halberstam--Richert~\cite[\S2.3, Thm.~2.2;
also \S6, Thm.~6.1]{HalberstamRichert1974},
Iwaniec--Kowalski~\cite[\S6.4, Prop.~6.7]{IwaniecKowalski2004}, or
Greaves~\cite[Ch.~4, eq.~(4.2)]{Greaves2001}) is: for even $r$,
\begin{equation}\label{eq:brun-upper}
  S(\mathcal{A}, \mathcal{P}_{z}, z)
  \;\leq\; \sum_{d \mid P(z), \, \omega(d) \leq r}
    \mu(d) \cdot \frac{X \cdot g_{V}(d)}{|d|}
    \;+\; \sum_{d \mid P(z), \, \omega(d) \leq r}
       |E_{V}(D, d, 0)|;
\end{equation}
and for odd $r$ the inequality reverses (lower bound). Here $g_{V}(d)
= \prod_{p \mid d} g_{V}(p) \cdot p$ for squarefree $d$ (extending
multiplicatively), and $|d| = d$ is the integer value of $d$. Taking
$r = 10$ even gives an upper bound; taking $r = 11$ odd gives a lower
bound. Since we want a lower bound, we take $r = 11$ odd in
\eqref{eq:brun-upper} (with the reversed inequality), and the
truncation error / loss factor at depth $r = 11$ is the same order as
at $r = 10$ (the bookkeeping is symmetric).

To estimate the gap between the truncated sum on the right of
\eqref{eq:brun-upper} and the untruncated product
$X \cdot W(z) = X \cdot \prod_{p} (1 - g_{V}(p))$, we use the
standard Bonferroni / Brun bound: the tail beyond truncation depth
$r$ is bounded by the $(r+1)$-st elementary symmetric function in the
$g_{V}(p)$, which is bounded above by
\begin{equation}\label{eq:brun-tail}
  \biggl| X \cdot W(z) - \sum_{d \mid P(z), \, \omega(d) \leq r}
    \mu(d) \cdot \frac{X \cdot g_{V}(d)}{d} \biggr|
  \;\leq\; X \cdot \frac{S_{1}^{r+1}}{(r+1)!} \cdot \exp(S_{1}),
\end{equation}
where $S_{1} := \sum_{p \in \mathcal{P}_{z}} g_{V}(p)$ and the
$\exp(S_{1})$ factor absorbs the lower-order combinatorial
multiplicities; this is the standard form recorded in
Halberstam--Richert~\cite[Lemma~2.2]{HalberstamRichert1974} or
Iwaniec--Kowalski~\cite[eq.~(6.7)]{IwaniecKowalski2004}. The
Stirling-style estimate $1/(r+1)! \leq (e/(r+1))^{r+1} \cdot
\sqrt{2\pi(r+1)}^{-1}$ gives the cleaner form
\begin{equation}\label{eq:brun-stirling}
  \biggl| \text{truncated sum} - X \cdot W(z) \biggr|
  \;\leq\; X \cdot W(z) \cdot \delta_{r},
  \qquad
  \delta_{r} \;\leq\; \frac{(e \cdot S_{1})^{r}}{r!}.
\end{equation}

\paragraph{Numerical bound on $S_{1}$.} By Remark~\ref{rem:loadbearing}
in \S\ref{sec:kappa}, $S_{1} = \sum_{p > 3, p \nmid m} g_{V}(p)$ is
bounded by an absolute constant; explicitly,
\begin{equation}\label{eq:S1-bound}
  S_{1} \;\leq\; \sum_{p > 3} \frac{1}{H_{p}^{2}}
  \;\leq\; \sum_{H = 2}^{\infty} \frac{M(H)}{H^{2}},
\end{equation}
where $M(H) := \#\{p : H_{p} = H\}$. Using the crude bound $M(H) \leq
H$ (which holds whenever $H \leq z$, via $H_{p} \mid p - 1$ giving at
most $H$ candidate primes $\leq H + 1$) gives the divergent
$\sum_{H} 1/H$, capped by $H \leq z$ to $S_{1} \leq C \log \log z$
unconditionally --- enough but not sharp. The
Erd\H{o}s--Pomerance / Pappalardi--Susa decay
(Lemma~\ref{lem:small-H-count}) replaces $M(H) \leq H$ by $M(H) \leq
C \cdot H / (\log H)^{B}$ for the small-$H$ range, giving
$S_{1} \leq C_{0}$ with $C_{0}$ absolute. We adopt the safe
conservative bound $S_{1} \leq 1$.

The loss factor at $r = 11$ (using Stirling
$11! = 3.99 \times 10^{7}$) is then
\begin{equation}\label{eq:delta-numeric}
  \delta_{11} \;\leq\; \frac{(e \cdot S_{1})^{11}}{11!}
  \;\leq\; \frac{e^{11}}{11!}
  \;=\; \frac{5.99 \times 10^{4}}{3.99 \times 10^{7}}
  \;<\; 1.51 \times 10^{-3}.
\end{equation}
With the tighter $S_{1} \leq 0.5$ (well within reach of the Pomerance
bookkeeping), the loss factor improves to
$\delta_{11} \leq (e/2)^{11} / 11! \leq 3.5 \times 10^{-5}$, but
\eqref{eq:delta-numeric} suffices for the constants chain.
We adopt the conservative choice $c_{\mathrm{Brun}} = 1 - 10^{-2}
= 0.99$ for the constants chain in \S\ref{sec:constants}; the tighter
value $1 - 1.51 \times 10^{-3}$ gives an even sharper $D_{0}$ if
desired.

\paragraph{Remainder bound.} For $r = 10$ (or $r = 11$ for the lower
bound), the truncated sum on the right of \eqref{eq:brun-upper} ranges
over squarefree $d = p_{1} p_{2} \cdots p_{j}$ with $j \leq 10$ and
each $p_{i} \leq z = D^{1/3}$. Hence $d \leq z^{10} = D^{10/3}$, which
greatly exceeds the BV-style level $Q = D^{1/2 - \varepsilon}$ for any
reasonable $\varepsilon$ --- so the cumulative remainder sum needs
splitting. We use the standard device of separating into ``small $d$''
($d \leq Q$, where Proposition~\ref{prop:BV} applies directly) and
``large $d$'' ($d > Q$, where the per-$d$ error is bounded trivially
by $|T_{D}|/d$ plus the standard perimeter term). The dominant
contribution is from small $d$: by Proposition~\ref{prop:BV},
\begin{equation}
  \sum_{d \leq Q, \, \mu^{2}(d) = 1, \, \omega(d) \leq 10}
    |E_{V}(D, d, 0)| \;\leq\; O(D^{2} / (\log D)^{A}).
\end{equation}
The large-$d$ contribution is $o(D^{2})$ by the trivial bound
$|E_{V}(D, d, 0)| \leq |T_{D}|/d + O(D \cdot H_{d})$ summed over
$d \in (Q, z^{10}]$ with $\omega(d) \leq 10$ (the count of such $d$
is bounded by $\binom{\pi(z)}{10} \cdot z^{-10\varepsilon}$ which is
$o(1)$ for $\varepsilon > 0$). Combining,
\eqref{eq:R10-bound} holds.

\paragraph{Conclusion.} The truncated Brun lower bound (depth $r = 11$,
reversed inequality in \eqref{eq:brun-upper}) gives
$S(\mathcal{A}, \mathcal{P}_{z}, z) \geq X \cdot W(z) \cdot
c_{\mathrm{Brun}} - R_{10}(z, Q)$ with $c_{\mathrm{Brun}} = 1 -
\delta_{11}$ and $\delta_{11} \leq 1.51 \times 10^{-3}$
\eqref{eq:delta-numeric}, well below the conservative
$10^{-2}$. The remainder $R_{10}(z, Q) = O(D^{2} / (\log D)^{A})$
follows from Proposition~\ref{prop:BV} as above. This establishes
\eqref{eq:brun-pure-lower}.
\end{proof}

\begin{remark}[Reproducibility checklist for \S\ref{sec:sieve}]
\label{rem:sieve-checklist}
For a referee re-deriving the Brun pure-sieve lower bound, the
load-bearing inputs are:
\begin{enumerate}
  \item The Brun pure-sieve upper/lower bound \eqref{eq:brun-upper}
    with the Bonferroni tail bound \eqref{eq:brun-tail}:
    \cite[\S2.3 Thm.~2.2 + Lemma~2.2]{HalberstamRichert1974} or
    \cite[\S6.4 Prop.~6.7 + eq.~(6.7)]{IwaniecKowalski2004}.
  \item The Stirling-form loss bound \eqref{eq:brun-stirling}: standard.
  \item A bound on $S_{1} = \sum_{p \in \mathcal{P}_{z}} g_{V}(p)$
    \eqref{eq:S1-bound}. The crude bound $M(H) \leq H$ via $H_{p}
    \mid p - 1$ gives the divergent $\sum 1/H$, but capped by
    $p \leq z$ this gives $S_{1} \leq C \log \log z$ unconditionally
    --- enough but not sharp. The Erd\H{o}s--Pomerance /
    Pappalardi--Susa decay (Lemma~\ref{lem:small-H-count}) replaces
    this by $S_{1} \leq C_{0}$ absolute. The conservative choice
    $S_{1} \leq 1$ is safe; tighter is better but not load-bearing.
  \item The BV bound (Proposition~\ref{prop:BV}) for the cumulative
    remainder.
\end{enumerate}
The combinatorial loss $\delta_{11} \leq 1.51 \times 10^{-3}$
\eqref{eq:delta-numeric} (or even smaller with the tight $S_{1}$) is
well below the conservative $\delta_{r} \leq 10^{-2}$ adopted in the
constants chain (\S\ref{sec:constants}).
\end{remark}

\paragraph{Conclusion of \S4.}
At $\kappa_{V} = 0$, Brun's truncated combinatorial sieve gives the
lower bound \eqref{eq:brun-pure-lower} with explicit
$c_{\mathrm{Brun}} = 1 - \delta_{11}$ and explicit $\delta_{11} \leq
1.51 \times 10^{-3}$. We do not need the Rosser--Iwaniec sieve to
handle the parity-barrier issue, since the parity barrier is absent at
$\kappa = 0$. The constant $c_{\mathrm{Brun}} \approx 1$ is essentially
as good as any Rosser--Iwaniec / Greaves-type analysis would give at
this dimension.

\paragraph{Net effect on the main theorem.}
Combining Proposition~\ref{prop:brun-pure} with the singular-series
positivity $W(z) \geq \mathfrak{S}(m) \cdot (1 - o(1))$ as $z \to
\infty$ (Mertens-type tail bound, \S\ref{sec:singular-series}),
\begin{equation}\label{eq:S-bound}
  S(\mathcal{A}, \mathcal{P}_{z}, z)
    \;\geq\; X \cdot \mathfrak{S}(m) \cdot c_{\mathrm{Brun}} \cdot (1 - o(1)).
\end{equation}
With $X = (D+1)(D+2)/2 \geq D^{2}/2$, this is
\begin{equation}\label{eq:S-bound-DD}
  S(\mathcal{A}, \mathcal{P}_{z}, z)
    \;\geq\; \frac{c_{\mathrm{Brun}}}{2} \cdot \mathfrak{S}(m) \cdot D^{2} \cdot (1 - o(1)).
\end{equation}

This is the count of $V(m, k, \ell)$ in $T_{D}$ free of prime factors
$\leq z = D^{1/3}$. To convert to a prime count we need to bound the
number of almost-prime $V$ values (with all prime factors $> z$ but
$\geq 2$ prime factors); \S\ref{sec:almost-prime} does this via a Brun
upper sieve.
